\uuid{TAT7}
\exo7id{4321}
\titre{exo7 4321}
\auteur{quercia}
\organisation{exo7}
\datecreate{2010-03-12}
\isIndication{false}
\isCorrection{true}
\chapitre{Intégration}
\sousChapitre{Intégrale de Riemann dépendant d'un paramètre}
\cpge{mpsi-mp2i.fonctions-valeurs-reelles-complexes.b, pcsi.techniques-calcul-fonctions.b, ptsi.techniques-calcul-fonctions.b, mpsi-mp2i.analyse-asymptotique.b, pcsi.analyse-asymptotique.b, ptsi.analyse-asymptotique.b, bcpst1.derivation.a, bcpst1.developpements-limites.a}
\module{Analyse}
\niveau{L2}
\difficulte{}

\contenu{
\question{
Soit $a > 0$. Donner le DL en $x=1$ à l'ordre 3 de
$f(x) =  \int_{t=a/x}^{ax} \frac{\ln t}{a^2+t^2}\,d t$.
}
\reponse{
$f'(x) = \frac{2\ln x}{a(1+x^2)}  \Rightarrow  f(1+h) = \frac1{2a}(h^2-h^3) +  o(h^3)$.
}
}
